\subsection{What is a Circle?} 
As a geometric figure, how do we define a circle?
\begin{definition}[Circle]
Given a point $O$ in the plane and a real number $r>0$, the \term{circle of radius $r$ with center $O$} is the set of all points $C$ with $d(O,C)=r$.
\end{definition}
So a circle consists of all points at some given distance from the center.
\begin{icpic}
	\fill (0,0) circle (0.1) node[anchor=north]{$O$};
	\draw[geom,blue] (0,0) circle (3);
	\draw[geom] (0,0) -- node[anchor=north]{$r$} (30:3);
	\fill (30:3) circle (0.1);
\end{icpic}
\begin{note}
The number $r$ is called \emph{the} radius of the circle. A segment running from the center to the circle is called \emph{a} radius. Each radius has length $r$.
\end{note}
We can use this definition to come up with an equation for a circle.
\begin{itemize}
\item We call the \makeblank[1in] $O=(h,k)$.
\item We're writing an equation for a shape in the plane, so we call the point on the circle $P=(x,y)$
\item Then $d(O,P)=\makeblank$
\item Our basic equation is $\makeblank=r$
\item But square roots are hard to work with, so we \makeblank[1in] both sides.
\item This gives us \makeblank
\end{itemize}
\begin{displaybox}[Equation of a Circle (center-radius form)]
The circle centered at $(h,k)$ with radius $r$ is the graph of the equation
\[
(x-h)^2+(y-k)^2=r^2
\]
\end{displaybox}